\section{Extensions}\label{sec:extensions}

The baseline treats Europe as one actor, gives it a single defection option, and lets the United States keep its chips flowing while the two sides bargain. Each extension relaxes one of these assumptions. The first asks who in Europe actually decides. The second gives Europe the power to hurt the United States during a dispute, which is where the patron's dependence on ASML returns. The third adds partial defection.

\subsection{Who decides: nested chokepoints and veto players}\label{sec:nested}

Three features of the real decision structure matter. The Netherlands licenses ASML's exports to China, but a licence permits a shipment; it does not order one. ASML itself decides whether to ship, and its own exposure to US law, customers, suppliers, and shareholders is large. Turning ASML against the United States, finally, is not a Dutch decision at all: exports to the United States move under an EU general authorization, and the instrument designed for retaliation, the Anti-Coercion Instrument (ACI), requires qualified-majority decisions in the Council.

Let the firm's China profit from exit be $y$ and its expected loss from US measures aimed at it be $f^F$. If $y<f^F$, the firm will not ship unless the state compels it at cost $k_s \geq 0$, for instance under emergency powers like those the Dutch government used against Nexperia in 2025. For the veto-player results, let each member state $i$ have its own exit premium $Q_i = \omega_i - c_i + \chi_{r,i} r_i$, computed from its own share of China's offer and its own exposure to US and Chinese instruments. For the ACI, let $q_i$ be member $i$'s net gain from adopting a countermeasure against the United States, and let $\mathcal{B}$ be the family of blocking coalitions: at least four members holding more than 35 percent of EU population.

\begin{proposition}[Nested chokepoints]\label{prop:nested}
\begin{enumerate}[label=(\alph*)]
\item \emph{Firm compliance.} Europe's effective exit premium is $Q - k_s\mathbb{1}\{y<f^F\}$. If the United States can raise $f^F$ above $y$ at negligible cost, Europe's leverage is at most $[Q-k_s]^+$.
\item \emph{Holder veto.} If exit requires the consent of member $h$ alone, the United States pays $[Q_h]^+$. A conditional benefit granted to $h$ alone and withdrawn upon exit raises $c_h$ one for one and lowers this payment.
\item \emph{Blocking minorities.} If no one bids for adoption, the United States defeats an ACI countermeasure at cost $K(0) = \min_{S\in\mathcal{B}} \sum_{i\in S}[q_i]^+$. If a rival bidder values adoption at $B>0$ and moves second, the United States must make every qualified majority cost at least $B$; its cost $K(B) = \min_{\tau \geq 0}\sum_i \tau_i$ subject to $\sum_{i\in S}[\tau_i - q_i]^+ \geq B$ for every qualified majority $S$ can exceed $K(0)$ by more than $B$ and generally requires payments to members outside any minimal blocking coalition \citep{groseclose1996buying}.
\end{enumerate}
\end{proposition}

Part (a) is the institutional point in \citet{farrell2019weaponized} turned around. A state can weaponize a network hub only if its domestic institutions give it control over the hub. ASML meets the network condition but sits inside networks the United States controls: US-origin subsystems, US persons, and US customers. Europe's chokepoint is nested in American ones, and the exit threat is only as credible as the Dutch state's willingness to compel its own champion to break US law. Part (b) is divide and rule. Granting the Netherlands, and only the Netherlands, preferential AI access, as the scenario's tier system does (the real January 2025 AI Diffusion Rule similarly split the EU, placing ten members in its top tier for reasons unrelated to ASML), gives the veto player more to lose from exit than its neighbours. Part (c) corrects a piece of arithmetic. Under the 2026 Council population figures \citep{council2025population}, the states the scenario says opposed or abstained on the 2029 countermeasure (the Netherlands, Ireland, Poland, the three Baltic states, and Italy) hold 28.05 percent of EU population, short of a blocking minority, and any blocking coalition that excludes Germany, France, Italy, and Spain needs at least twelve members. The pivotal members on a measure against the United States are large states, not the eastern flank.

\subsection{Voice through hold-up: time to pain}\label{sec:holdup}

The baseline assumes that bargaining takes place under the status quo, with ASML servicing US-aligned fabs throughout. The scenario's 2031 standoff is different. Once Europe refuses the ultimatum, Washington cuts off frontier AI, and Europe could interrupt servicing, spare parts, and tool deliveries in return. The scenario claims that Washington ``has amassed enough EUV machines and maintenance know-how to forego ASML access for longer than Europe can forego access to frontier AI''. That is a claim about time to pain, and the patron's dependence on the chokepoint enters through it.

Let Europe and the United States dispute a divisible issue: the scope of US control over ASML, or the terms of Europe's AI access. An agreement leaves Europe a share $p \in [0,1]$, with payoffs $-(1-p)s_E$ to Europe and $(1-p)s_U$ to the United States. Until they agree, Europe bears a flow cost $c_E$ from date $\tau_E$ on, and the United States bears $c_U$ from date $\tau_U$ on. The buffers are what each side can live on before the cutoff hurts. For the United States, that means stockpiled tools, spare parts, and transferred maintenance know-how. For Europe, it means compute on European soil and stored model weights. Offers arrive at intervals $\Delta$; the United States proposes with probability $\psi$ and Europe otherwise. We take $\Delta \to 0$ and assume Europe's exit option is unattractive. Write $\hat c_i \equiv c_i/s_i$ and $D \equiv \tau_U - \tau_E$.

\begin{proposition}[Time to pain]\label{prop:holdup}
Suppose $\psi \hat c_E \neq (1-\psi)\hat c_U$.
\begin{enumerate}[label=(\alph*)]
\item The equilibrium is unique and agreement is immediate. If Europe is the less vulnerable side once both buffers are exhausted, $\psi\hat c_E < (1-\psi)\hat c_U$, its share is $p^*=1$ when $D\leq 0$ and $p^* = [1-\psi \hat c_E D]^+$ when $D>0$. If instead $\psi\hat c_E > (1-\psi)\hat c_U$, then $p^*=0$ when $D\geq 0$ and $p^* = \min\{1, (1-\psi)\hat c_U |D|\}$ when $D<0$.
\item Let $\tau_E = 0$ and suppose the United States bears no cost unless Europe interrupts servicing. Interruption, which raises Europe's flow cost to $c_E+\ell$ and imposes $c_U$ on the United States after $\tau_U$, is played in a Markov perfect equilibrium if and only if $(1-\psi)\hat c_U > \psi (c_E+\ell)/s_E$ and $\tau_U < H_E \equiv s_E/[\psi(c_E+\ell)]$. Europe's share is then $1 - \tau_U/H_E$.
\end{enumerate}
\end{proposition}

The proposition separates the two quantities \citet{keohane1977power} called sensitivity and vulnerability. The buffer gap $D$ measures sensitivity: who feels the cutoff first. The post-buffer ranking of normalized costs measures vulnerability: who could hold out longer once neither has a cushion. Hold-up gives Europe leverage only if it wins on vulnerability and the US buffer is shorter than Europe's endurance horizon $H_E$. In the interior, each unit of US buffer costs Europe $\psi c_E$, Europe's own rate of pain weighted by US agenda power; the US cost $c_U$ matters only for who wins on vulnerability. This is where $V$ re-enters the analysis. The more US frontier compute depends on continued ASML servicing, the larger $c_U$. The scenario's claim that a stockpile lets Washington outlast Europe amounts to the claim that $\tau_U \geq H_E$. Its Epilogue's call for compute on European soil amounts to raising $\tau_E$. Servicing and spare parts matter more than tools here: installed EUV machines keep running, but they need engineers and consumable parts, so the US buffer for servicing is short even if its stock of tools is large.

\subsection{Hedging and marginal deterrence}\label{sec:hedge}

The scenario's 2028 crisis, like the real disputes of 2023--2025, was not about exit. It was about whether ASML could keep selling and servicing older immersion tools in China. Add a third alignment, hedging ($H$), with China's credible offer $\omega_H < \omega$ and exit costs $c_H$, and suppose the United States, after observing a hedge, may withhold an instrument $k$ that also deters exit. Let $D_k$ be the instrument's contribution to Europe's exit costs ($\kappa S$ for security, $\alpha$ for AI) and $g_k(H)$ the United States' ex post gain from withholding it from a hedging Europe.

\begin{proposition}[Marginal deterrence]\label{prop:hedge}
Withholding instrument $k$ to punish hedging is credible if and only if
\begin{equation}
m_k + g_k(H) \;\geq\; [Q + D_k]^+ - [Q]^+ .
\end{equation}
Enforcing jurisdiction over the hedged flows themselves leaves $Q$ unchanged and is credible whenever its direct cost is small. If $Q + D_k > 0$ and audience costs are small, the United States cannot use security or AI linkage against hedging, and must rely on flow-specific jurisdiction.
\end{proposition}

A patron cannot spend the sanctions that deter full defection on partial defection without weakening deterrence of the full version \citep{stigler1970optimum, shavell1992note}. The prediction is that Washington polices legacy-tool sales with targeted jurisdiction over those sales (de minimis rules for lithography tools bound for China, the foreign direct product rule, servicing restrictions) and keeps security and AI linkage in reserve. This matches the scenario's 2028 crisis, where the threat is the FDPR, and it matches the actual US sequence from 2019 to 2024. Jurisdictional threats of this kind also work under uncertainty about US resolve. If the firm-level penalty $f_H$ is existential for ASML while the revenue at stake is a fraction of its sales, even a small probability that Washington means it is enough. That is the formal content of the scenario's line that ASML ``cannot afford to call Washington's bluff''.
